\begin{abstract}

Large text embedding models are central to semantic search, clustering, and retrieval-augmented generation (RAG), where quality depends on neighborhood structure over an indexed corpus. Existing embedding distillation methods transfer local teacher signals, from embeddings and token states to batch-level similarities. However, in retrieval and RAG, an embedding model is judged after it becomes an index: small distortions in corpus-level neighborhoods can change which evidence is surfaced, even when individual student vectors remain close to the teacher. This exposes a mismatch between local distillation losses and the neighborhood dynamics retrieval depends on. The key object to preserve is therefore not a set of vectors, but the teacher's retrieval operator over the corpus. We propose \textbf{HeatGeo}, a corpus-level geometry distillation framework that models this behavior as a semantic random walk on a teacher-induced neighborhood graph. HeatGeo caches teacher embeddings, builds a mutual $k$-nearest-neighbor graph, and derives multi-scale heat-kernel diffusion distributions describing how semantic mass spreads from each example to nearby regions and higher-order communities. The student matches these distributions over sampled candidate neighborhoods, while a spectral manifold objective anchors global topology using graph Laplacian eigenvectors. We evaluate on STS, text classification, and pair classification to test whether corpus-level dynamics preserve retrieval geometry better than pointwise, token-level, and mini-batch relational baselines.

\end{abstract}
